\section{Related work}
\label{sec:related}

Social intersects crisis informatics, news-event databases, automated claim
verification, retrieval-supported generation, and digital investigation practice.
These traditions address different units: messages, reported events, claims,
answers, or preserved media. Their outputs therefore require different validation.
A retrieved article may be relevant without supporting a claim; several reports
may share one origin; and a decoded video may remain unverified as to location,
time, or interpretation. We use this distinction to position the inspected
software, rather than treating adjacent systems as interchangeable competitors.
Table~\ref{tab:related} summarizes their documented purposes and comparison limits.

Crisis informatics provides the immediate problem setting. Imran et
al.~\cite{imran2015crisis} organize emergency-message processing around filtering,
classification, ranking, aggregation, extraction, and summarization. This framing
connects heterogeneous collection with the analyst's need to reduce information
overload. Social adopts a similar workflow perspective but studies a bounded
incident reconstruction and its inspectable artifacts. Its platform panels and
case-specific selection rules do not constitute a representative population
sample. Successful execution of a filter or collector consequently cannot
establish completeness, situational accuracy, or usefulness to emergency
responders without a separate evaluation of those outcomes.

News-event infrastructures show that organizing observations into events is an
established task. GDELT separates CAMEO-coded event records from news mentions and
provides extraction confidence and geographic attributes; its GKG supplies
document-level entities, themes, and contextual
metadata~\cite{gdelt2015events,gdelt2015gkg}. ACLED uses researcher coding and a
defined political-disorder taxonomy, with explicit date and location precision
and source information~\cite{acledCodebook}. Event Registry groups multilingual
news articles, extracts event attributes, and supports event search and
exploration~\cite{leban2014eventregistry}. These works motivate explicit event
units and distinctions between reports and occurrences. Social does not claim
novelty for extraction or clustering. Its narrower research question concerns
how capture receipts, evidence passages, claim assessments, and event hypotheses
can remain inspectably connected within a bounded investigation.

Automated fact-checking introduces a complementary requirement: evaluate the
evidence as well as the predicted judgment. FEVER combines claims derived from
Wikipedia with evidence selection and supported, refuted, or insufficient-information
labels~\cite{thorne2018fever}. AVeriTeC addresses real-world claims using web
evidence, question--answer decompositions, and justifications, explicitly attending
to temporal leakage and evidence sufficiency~\cite{schlichtkrull2023averitec}.
InFact implements a staged, retrieval-based LLM fact-checking
pipeline~\cite{rothermel2024infact}. These are relevant task formulations and
evaluation precedents. Social's contract checks and incident regressions do not
amount to a FEVER or AVeriTeC benchmark result, and no accuracy ranking against
these systems is implied. Claim-level correctness and reconstruction quality
remain separate from software integration success.

Retrieval-augmented generation combines learned generation with external
retrieved memory~\cite{lewis2020rag}. ALCE makes citation quality a distinct
evaluation dimension alongside correctness and
fluency~\cite{gao2023alce}. This distinction matters for generated investigation
reports: a syntactically present citation is not necessarily relevant, sufficient,
or faithful to the associated sentence. Social's source and passage references
provide objects that can be inspected, while their existence alone supplies no
guarantee of semantic support. Neither retrieval relevance nor a model's verbal
confidence should be interpreted as an independently calibrated probability that
an incident claim is true.

Diffusion observatories and media-verification tools address further parts of
the workflow. Hoaxy tracks the circulation of URLs from monitored sources and
their fact-checking counterparts; its 2016 description explicitly distinguishes
tracking from performing fact checking~\cite{shao2016hoaxy}. InVID and WeVerify
combine keyframe extraction, reverse-image search, metadata inspection, forensic
assistance, and guided analyst tools~\cite{invidVerificationPlugin,weverifyVerificationPlugin}.
They establish substantial prior art for integrated verification assistance.
Their documented functions do not make every diagnostic result a factual
verdict. Likewise, Social's local media indexing and candidate transcripts must
be distinguished from independent authentication and source-context verification.

The Berkeley Protocol separates collection, preservation, verification,
investigative analysis, and reporting within a professional and ethical
investigation process~\cite{berkeleyProtocol2022}. It is methodological guidance,
not a software certification. Research-software practices concerning organized
data, documented transformations, version control, and reproducible workflows
provide a complementary engineering frame~\cite{wilson2017practices}. Social's
contribution is assessed against concrete implemented artifacts and bounded
checks within these established traditions. Neither adopting a terminology nor
passing local tests establishes protocol compliance, evidentiary admissibility,
universal platform coverage, or semantic validity across incidents. Those broader
claims require additional evidence and an evaluation design matched to them.

\begin{table*}[t]
\centering
\footnotesize
\setlength{\tabcolsep}{3pt}
\renewcommand{\arraystretch}{1.12}
\caption{Qualitative positioning from the cited primary sources and the present
code audit. Cells describe documented functions or validation scope, not measured
superiority. Historical papers and tool pages do not establish current platform
access.}
\label{tab:related}
\begin{tabularx}{\textwidth}{@{}p{20mm}YYYYY@{}}
\toprule
System & Primary purpose & Provenance representation & Event reasoning & Uncertainty & Evaluation / comparison \\
\midrule
GDELT & News-event and contextual metadata & Event/mention identifiers; document links & CAMEO actions; contextual fields & Extraction confidence; geographic estimates & Codebook describes fields; no shared incident benchmark here \\
ACLED & Coded political-disorder events & Sources and event notes & Researcher-coded event taxonomy & Date/location precision; coding conventions & Methodology and review process; different observation unit \\
Event Registry & News-to-event organization & Associated news articles & Multilingual clustering and attributes & Language and clustering scope & 2014 system description; no common test here \\
Hoaxy & URL-sharing observatory & Monitored URLs and sharing observations & Circulation analysis & Selected sources and platform sample & Preliminary sharing analysis; no truth adjudication \\
InVID--WeVerify & Analyst media-verification assistance & Context, metadata and search results & Analyst-guided media examination & Diagnostic tools require interpretation & Official tool descriptions; no joint performance test \\
Social & Bounded incident evidence workflow & Capture hashes, paths, passages and typed links & Claim assessments; event hypotheses & Access states and unresolved evidence & Local regression/integrity checks; limited semantic validation \\
\bottomrule
\end{tabularx}
\end{table*}
